\chapter{Nail Clubbing}

\section*{Definition}
Nail clubbing (Hippocratic fingers) is a bulbous enlargement of the distal fingers and toes with loss of the normal Lovibond angle (\ensuremath{>}180\textdegree{}), nail bed fluctuation, and increased hyponychial angle. It may be idiopathic, familial, or a sign of underlying cardiopulmonary, hepatic, or neoplastic disease.

\section*{What Happens in Your Body}
Clubbing results from platelet-derived growth factor (PDGF) and vascular endothelial growth factor (VEGF) release when megakaryocytes and platelet precursors become trapped in the distal digital microvasculature rather than undergoing fragmentation in the pulmonary circulation. In cyanotic heart disease or intrapulmonary shunting, megakaryocytes bypass pulmonary filtration and reach the fingertips, releasing growth factors that induce edema, hyperemia, and connective tissue hyperplasia of the nail bed.

\section*{Causes}
\begin{itemize}
  \item \textbf{Pulmonary:} Lung cancer (most common malignant cause, 80\% of clubbed people), bronchiectasis, empyema, pulmonary fibrosis, cystic fibrosis, tuberculosis, lung abscess
  \item \textbf{Cardiac:} Cyanotic congenital heart disease (Tetralogy of Fallot, transposition), infective endocarditis, atrial myxoma
  \item \textbf{Hepatic:} Cirrhosis, hepatocellular carcinoma, biliary atresia
  \item \textbf{Gastrointestinal:} Inflammatory bowel disease (Crohn's disease, ulcerative colitis), celiac disease, intestinal lymphoma
  \item \textbf{Other:} Thyroid acropachy, unilateral clubbing (Pancoast tumor, axillary aneurysm), familial (autosomal dominant, benign)
\end{itemize}

\section*{When to See a Doctor}
See your doctor if:
\begin{itemize}
  \item New-onset clubbing in an adult, especially a smoker (must rule out lung cancer)
  \item Clubbing with cough, hemoptysis, dyspnea, or chest pain
  \item Associated cyanosis, clubbing, and digital edema in congenital heart disease
  \item Unilateral clubbing or finger swelling with Horner syndrome (suspect Pancoast tumor)
\end{itemize}

\section*{Related Symptoms}
\begin{itemize}
  \item Hypertrophic pulmonary osteoarthropathy (HPOA), cyanosis, arthralgia, digital edema, periostitis
\end{itemize}
\textbf{Important:} Digital clubbing with HPOA (periostitis, arthralgia, clubbing) is a paraneoplastic syndrome most commonly associated with non-small cell lung cancer and may improve with tumor resection.

\section*{Self-Care / Home Management}
\begin{itemize}
  \item No self-care reverses clubbing; management focuses on identifying and treating the underlying cause
  \item Smoking cessation if indicated
  \item No intervention is needed for idiopathic or familial clubbing
\end{itemize}

\section*{How Your Doctor Will Evaluate You}
\begin{itemize}
  \item Physical examination: Schamroth method (loss of diamond-shaped window placing dorsal fingers together), profile sign, hyponychial angle \ensuremath{>}192\textdegree{}, phalangeal depth ratio \ensuremath{>}1.0
  \item Chest X-ray as first-line imaging; CT chest if suspicious findings or unexplained clubbing
  \item Pulmonary function tests, arterial blood gas, echocardiography, or abdominal imaging as guided by clinical presentation
\end{itemize}

\section*{Treatment Options}
\begin{itemize}
  \item Treat underlying condition: lung cancer resection, antibiotics for empyema/bronchiectasis, surgical repair of congenital heart disease
  \item Inflammatory bowel disease: Disease-modifying agents (biologics, immunomodulators)
  \item Clubbing typically resolves after successful treatment of the underlying cause (e.g., tumor resection, infection clearance)
  \item No specific treatment is needed for benign familial clubbing
\end{itemize}

\section*{References}
\begin{thebibliography}{99}
\bibitem{ref1} Spicknall KE, Zirwas MJ, English JC III. "Clubbing: an update on diagnosis, differential diagnosis, and clinical relevance." J Am Acad Dermatol. 2005;52(6):1020-1028.
\bibitem{ref2} Myers KA, Farquhar DR. "Does this patient have clubbing?" JAMA. 2001;286(3):341-347.
\bibitem{ref3} Dickinson JP. "The natural history of clubbing." Thorax. 2008;63(6):495-498.
\bibitem{ref4} Fauci AS, Braunwald E, Kasper DL, et al. "Digital clubbing." In: Harrison's Principles of Internal Medicine. 21st ed. McGraw-Hill; 2022.
\bibitem{ref5} Sarkar M, Mahesh DM, Madabhavi I. "Digital clubbing." Lung India. 2012;29(4):354-362.
\bibitem{ref6} West JB. "Digital clubbing." UpToDate. Wolters Kluwer; 2025.
\end{thebibliography}

\clearpage
